\section{Django Framework and MVT Architecture}

The AI-Based Trip Planner Web Application is developed using the \textbf{Django} framework \cite{django2026}, one of the most widely used Python-based web development frameworks. Django follows the \textbf{Model-View-Template (MVT)} architectural pattern, which separates the application into independent layers for data handling, business logic, and presentation rendering.

\begin{figure}[H]
    \centering
    \begin{tikzpicture}[
        node distance=1.2cm,
        box/.style={draw, rounded corners, minimum width=4cm, minimum height=0.85cm, align=center, font=\small},
        arrow/.style={-Stealth, thick}
    ]
        \node[box, fill=blue!15] (req) {HTTP Request (User)};
        \node[box, fill=orange!15, below=of req] (url) {URL Router (urls.py)};
        \node[box, fill=green!15, below=of url] (view) {View (views.py)};
        \node[box, fill=purple!15, below left=1.2cm and 1.5cm of view] (model) {Model (models.py)};
        \node[box, fill=red!15, below right=1.2cm and 1.5cm of view] (template) {Template (.html)};
        \node[box, fill=gray!15, below=of model] (db) {Database (SQLite3)};
        \node[box, fill=blue!10, below=of template] (resp) {HTTP Response (Browser)};

        \draw[arrow] (req) -- (url);
        \draw[arrow] (url) -- (view);
        \draw[arrow] (view.south west) -- (model.north);
        \draw[arrow] (view.south east) -- (template.north);
        \draw[arrow] (model) -- (db);
        \draw[arrow] (model.east) -- ++(0.5,0) |- (template.west);
        \draw[arrow] (template) -- (resp);
    \end{tikzpicture}
    \caption{Django MVT Architecture}
    \label{fig:mvt-architecture}
\end{figure}

The MVT architecture improves code organization, scalability, and maintainability by enforcing separation of concerns:

\begin{itemize}
    \item \textbf{Model}: Defines the database schema and handles data persistence through ORM queries.
    \item \textbf{View}: Contains the application business logic, processes HTTP requests, and returns responses.
    \item \textbf{Template}: Renders the HTML presentation layer with dynamic data injected by views.
\end{itemize}

\subsection{Models Layer}

The \texttt{Trip} model defines the database schema for storing trip records:

\begin{lstlisting}[language=Python, caption={Trip Model Definition}, label={lst:trip-model}]
from django.db import models
from django.utils import timezone

class Trip(models.Model):
    CLIMATE_CHOICES = [
        ('hot', 'Hot'),
        ('cold', 'Cold'),
        ('moderate', 'Moderate'),
    ]
    TRAVEL_TYPE_CHOICES = [
        ('adventure', 'Adventure'),
        ('relaxation', 'Relaxation'),
        ('cultural', 'Cultural'),
        ('family', 'Family'),
    ]
    destination = models.CharField(max_length=255, default='Unknown')
    budget = models.IntegerField(default=0)
    days = models.IntegerField(default=1)
    climate = models.CharField(max_length=20, choices=CLIMATE_CHOICES)
    travel_type = models.CharField(max_length=20, choices=TRAVEL_TYPE_CHOICES)
    itinerary = models.TextField(blank=True)
    plan_data = models.JSONField(default=dict, blank=True)
    image = models.ImageField(upload_to='trip_images/', blank=True, null=True)
    is_deleted = models.BooleanField(default=False)
    deleted_at = models.DateTimeField(null=True, blank=True)
    created_at = models.DateTimeField(auto_now_add=True)
\end{lstlisting}

\subsection{Views Layer}

The views layer handles business logic and request processing. The \textbf{request handling workflow} operates as follows:

\begin{figure}[H]
    \centering
    \begin{tikzpicture}[
        node distance=0.9cm,
        box/.style={draw, rounded corners, minimum width=6.5cm, minimum height=0.8cm, align=center, font=\small},
        arrow/.style={-Stealth, thick}
    ]
        \node[box, fill=blue!10] (r1) {Receive HTTP Request};
        \node[box, fill=blue!10, below=of r1] (r2) {Authenticate user session};
        \node[box, fill=green!10, below=of r2] (r3) {Parse and validate request data};
        \node[box, fill=green!10, below=of r3] (r4) {Call ML service for prediction};
        \node[box, fill=orange!10, below=of r4] (r5) {Generate itinerary plan};
        \node[box, fill=orange!10, below=of r5] (r6) {Save trip record to database};
        \node[box, fill=purple!10, below=of r6] (r7) {Return JSON/Template response};

        \draw[arrow] (r1) -- (r2);
        \draw[arrow] (r2) -- (r3);
        \draw[arrow] (r3) -- (r4);
        \draw[arrow] (r4) -- (r5);
        \draw[arrow] (r5) -- (r6);
        \draw[arrow] (r6) -- (r7);
    \end{tikzpicture}
    \caption{Request Handling Workflow}
    \label{fig:request-workflow}
\end{figure}

\subsection{Templates Layer}

Django templates render dynamic HTML pages using template tags and variables. The system uses a base template (\texttt{base.html}) with \texttt{\{\% extends \%\}} inheritance for consistent layout across all pages.

\section{Authentication System}

The application implements Django's built-in authentication framework with custom views for secure user registration, login, and session management.

The authentication workflow is as follows:

\begin{figure}[H]
    \centering
    \begin{tikzpicture}[
        node distance=1.2cm,
        box/.style={draw, rounded corners, minimum width=6cm, minimum height=0.85cm,
                    align=center, font=\small},
        decision/.style={draw, diamond, aspect=2.5, minimum width=5cm,
                         align=center, font=\small},
        arrow/.style={-Stealth, thick}
    ]
        % Top: form submission
        \node[box, fill=blue!10] (s1) {User submits login / register form};

        % Decision diamond
        \node[decision, fill=yellow!25, below=1.2cm of s1] (d1) {Valid credentials?};

        % YES path — straight down
        \node[box, fill=green!15, below=1.2cm of d1] (s2)
              {Create session and redirect to dashboard};

        % NO path — to the left
        \node[box, fill=red!15, left=2.5cm of d1] (s3) {Display error message};

        % Arrows
        \draw[arrow] (s1)   -- (d1);
        \draw[arrow] (d1)   -- node[right, font=\scriptsize]{Yes} (s2);
        \draw[arrow] (d1.west) -- node[above, font=\scriptsize]{No} (s3.east);
    \end{tikzpicture}
    \caption{Authentication Workflow}
    \label{fig:auth-workflow}
\end{figure}

\subsection{User Registration}

The registration view accepts a username, email, and password. Passwords are validated for strength and hashed using Django's default PBKDF2 algorithm before storage. After successful registration, users are automatically redirected to the login page.

\section{Itinerary Generation Engine}

The itinerary generation system produces dynamic day-wise travel schedules based on the predicted destination and user-specified trip duration. Each destination in the system has a predefined activity database containing:

\begin{itemize}
    \item Daily sightseeing activities
    \item Recommended dining and local experiences
    \item Suggested timings for each activity
    \item Local travel recommendations
\end{itemize}

The itinerary engine automatically:
\begin{itemize}
    \item Truncates schedules for shorter trips (fewer days than the full itinerary)
    \item Extends schedules for longer trips by repeating activity cycles
    \item Rearranges activities dynamically to match trip duration
\end{itemize}

\begin{figure}[H]
    \centering
    \begin{tikzpicture}[
        node distance=1.0cm,
        box/.style={draw, rounded corners, minimum width=6.5cm, minimum height=0.85cm, align=center, font=\small},
        arrow/.style={-Stealth, thick}
    ]
        \node[box, fill=blue!15] (pred) {Predicted Destination};
        \node[box, fill=green!15, below=of pred] (lookup) {Lookup Destination Activity Database};
        \node[box, fill=orange!15, below=of lookup] (adjust) {Adjust Schedule to Trip Duration};
        \node[box, fill=purple!15, below=of adjust] (cost) {Calculate Cost Breakdown};
        \node[box, fill=red!15, below=of cost] (json) {Format as JSON Plan Data};
        \node[box, fill=gray!15, below=of json] (save) {Save to Trip Model};

        \draw[arrow] (pred) -- (lookup);
        \draw[arrow] (lookup) -- (adjust);
        \draw[arrow] (adjust) -- (cost);
        \draw[arrow] (cost) -- (json);
        \draw[arrow] (json) -- (save);
    \end{tikzpicture}
    \caption{Itinerary Generation Pipeline}
    \label{fig:itinerary-pipeline}
\end{figure}

\subsection{Cost Estimation Logic}

The cost estimation module distributes the user's total budget across standard travel expense categories using fixed percentage allocations:

\begin{itemize}
    \item \textbf{Accommodation}: 40\%
    \item \textbf{Food}: 25\%
    \item \textbf{Transportation}: 20\%
    \item \textbf{Activities}: 15\%
\end{itemize}

This breakdown totals exactly 100\%, ensuring accurate financial planning in the generated itinerary.

\section{API Integration and Django REST Framework}

The application uses \textbf{Django REST Framework} for API development and JSON serialization. REST APIs enable communication between frontend interfaces and backend services while supporting scalable application design.

\begin{figure}[H]
    \centering
    \begin{tikzpicture}[
        node distance=1.0cm,
        box/.style={draw, rounded corners, minimum width=5cm, minimum height=0.85cm, align=center, font=\small},
        arrow/.style={-Stealth, thick}
    ]
        \node[box, fill=blue!15] (client) {Frontend Client (JavaScript fetch)};
        \node[box, fill=green!15, below=of client] (api) {/api/recommendation/ (POST)};
        \node[box, fill=orange!15, below=of api] (serial) {TripRequestSerializer (Validation)};
        \node[box, fill=purple!15, below=of serial] (ml) {ML Prediction Service};
        \node[box, fill=red!15, below=of ml] (resp) {JSON Response (destination + itinerary)};

        \draw[arrow] (client) -- (api);
        \draw[arrow] (api) -- (serial);
        \draw[arrow] (serial) -- (ml);
        \draw[arrow] (ml) -- (resp);
        \draw[arrow, dashed] (resp.east) -- ++(1,0) |- (client.east);
    \end{tikzpicture}
    \caption{API Communication Architecture}
    \label{fig:api-architecture}
\end{figure}

\subsection{Request Validation}

The application uses \texttt{TripRequestSerializer} to validate incoming request data. Validation checks include:

\begin{itemize}
    \item Required fields presence (budget, days, climate, travel\_type)
    \item Numeric constraints (budget $> 0$, days between 1 and 30)
    \item Allowed climate values (hot, cold, moderate)
    \item Travel category validation (adventure, relaxation, cultural, family)
\end{itemize}

\subsection{JSON Response Structure}

The API returns structured JSON responses containing the predicted destination, budget breakdown, estimated expenses, daily itinerary, and travel metadata. An example response:

\begin{lstlisting}[language=Python, caption={API JSON Response Example}, label={lst:api-response}]
{
    "destination": "Manali",
    "budget": 25000,
    "days": 5,
    "itinerary": {
        "Day 1": "Arrival and local sightseeing",
        "Day 2": "Adventure sports and trekking",
        "Day 3": "Cultural exploration and shopping"
    },
    "cost_breakdown": {
        "accommodation": 10000,
        "food": 6250,
        "transportation": 5000,
        "activities": 3750
    }
}
\end{lstlisting}

\begin{table}[H]
    \centering
    \caption{API Endpoint Specifications}
    \label{tab:api-endpoints}
    \begin{tabularx}{\textwidth}{@{}lllX@{}}
        \toprule
        \textbf{Endpoint} & \textbf{Method} & \textbf{Payload} & \textbf{Response} \\
        \midrule
        \texttt{/api/recommendation/} & POST & \{budget, days, climate, travel\_type\} & JSON destination, itinerary, and cost breakdown \\
        \texttt{/} & GET & --- & Home page template \\
        \texttt{/plan-trip/} & GET & --- & Trip planning form template \\
        \texttt{/results/} & GET & data query param & Results view template \\
        \texttt{/my-trips/} & GET & --- & Dashboard template \\
        \texttt{/trips/<id>/delete/} & POST & CSRF token & Redirect to dashboard \\
        \bottomrule
    \end{tabularx}
\end{table}
